\documentclass[border=12pt]{standalone}
\usepackage[no-math]{fontspec}
\usepackage[brazil]{babel}
\usepackage{amsmath,amssymb}
\IfFileExists{C:/Windows/Fonts/Lato_Regular.ttf}{%
  \setmainfont{Lato}[
    Path = C:/Windows/Fonts/,
    UprightFont = Lato_Regular.ttf,
    BoldFont = Lato_Bold.ttf,
    ItalicFont = Lato_Italic.ttf,
    BoldItalicFont = Lato_BoldItalic.ttf
  ]
  \setsansfont{Lato}[
    Path = C:/Windows/Fonts/,
    UprightFont = Lato_Regular.ttf,
    BoldFont = Lato_Bold.ttf,
    ItalicFont = Lato_Italic.ttf,
    BoldItalicFont = Lato_BoldItalic.ttf
  ]
}{%
  \RequirePackage[default]{lato}
}
\renewcommand{\familydefault}{\sfdefault}

\usepackage{xcolor}
\usepackage{tikz}
\usepackage{pgfplots}
\pgfplotsset{compat=1.18}
\usetikzlibrary{arrows.meta,positioning,shapes.geometric,fit,backgrounds,calc,decorations.pathreplacing,patterns}

% Paleta institucional UEFS / Extensão Rural
\definecolor{erNavy}{HTML}{1F4268}
\definecolor{erNavyDark}{HTML}{152E4A}
\definecolor{erGold}{HTML}{B38A3D}
\definecolor{erGoldDark}{HTML}{8C6927}
\definecolor{erTeal}{HTML}{0E7C7B}
\definecolor{erForest}{HTML}{2D6A4F}
\definecolor{erAlert}{HTML}{9E2A2B}
\definecolor{erInk}{HTML}{26313D}
\definecolor{erMuted}{HTML}{5B6573}

\definecolor{bgNavy}{HTML}{EDF2F7}
\definecolor{bgGold}{HTML}{F8F3E9}
\definecolor{bgTeal}{HTML}{EBF6F5}
\definecolor{bgForest}{HTML}{EEF5F0}
\definecolor{bgAlert}{HTML}{FDF2F2}

\tikzset{
  >=Stealth,
  erBox/.style={
    draw=#1, fill=white, rounded corners=5pt, line width=1.15pt,
    inner sep=7pt, align=center, font=\footnotesize\sffamily, text=erInk
  },
  erTint/.style={
    draw=#1, fill=#1!8!white, rounded corners=5pt, line width=1.15pt,
    inner sep=7pt, align=center, font=\footnotesize\sffamily, text=erInk
  },
  erBadge/.style={
    fill=#1, text=white, rounded corners=9pt, inner xsep=8pt, inner ysep=3pt,
    font=\scriptsize\bfseries\sffamily
  },
  erTitle/.style={font=\normalsize\bfseries\sffamily, text=erNavy, align=center},
  erSub/.style={font=\footnotesize\sffamily, text=erMuted, align=center},
  erArrow/.style={->, >=Stealth, line width=1.35pt},
  erDArrow/.style={<->, >=Stealth, line width=1.35pt}
}

\hyphenpenalty=10000
\exhyphenpenalty=10000

\begin{document}
\begin{tikzpicture}[font=\sffamily]

% Título principal
\node[erTitle, font=\large\bfseries\sffamily] (title) at (0, 4.45)
  {Objetos de Mediação Sociotécnica na Comunicação Rural};

% Polos laterais de conhecimento
\node[
  draw=erNavy, fill=bgNavy, rounded corners=6pt, line width=1.25pt,
  text width=4.05cm, minimum height=1.85cm, inner sep=10pt, align=center
] (tech) at (-5.75, 2.35) {%
  {\color{erNavy}\small\bfseries Saber Técnico /\\Agronomia}\\[4pt]
  {\color{erMuted}\scriptsize Conhecimento científico\\e normativo}
};

\node[
  draw=erForest, fill=bgForest, rounded corners=6pt, line width=1.25pt,
  text width=4.05cm, minimum height=1.85cm, inner sep=10pt, align=center
] (farm) at (5.75, 2.35) {%
  {\color{erForest}\small\bfseries Saber Prático /\\Camponês}\\[4pt]
  {\color{erMuted}\scriptsize Experiência territorial\\e saber tradicional}
};

% Hub central: Objeto de Mediação
\node[
  draw=erGoldDark, fill=bgGold, rounded corners=8pt, line width=1.6pt,
  text width=3.9cm, minimum height=1.95cm, inner sep=10pt, align=center
] (hub) at (0, 2.35) {%
  {\color{erGoldDark}\normalsize\bfseries Objeto de\\Mediação}\\[4pt]
  {\color{erInk}\scriptsize Tradução horizontal\\sem hierarquia}
};

% Setas bidirecionais laterais
\draw[erDArrow, color=erNavy]   (tech.east) -- (hub.west);
\draw[erDArrow, color=erForest] (hub.east)  -- (farm.west);

% 4 Pílulas de exemplo abaixo (com espaçamento horizontal limpo)
\node[
  draw=erTeal, fill=bgTeal, rounded corners=12pt, line width=1.15pt,
  text width=3.4cm, minimum height=1.18cm, inner sep=7pt, align=center,
  font=\footnotesize\bfseries\sffamily, text=erTeal
] (p1) at (-6.45, -1.1) {Croquis \& Mapas\\Participativos};

\node[
  draw=erNavy, fill=bgNavy, rounded corners=12pt, line width=1.15pt,
  text width=3.4cm, minimum height=1.18cm, inner sep=7pt, align=center,
  font=\footnotesize\bfseries\sffamily, text=erNavy
] (p2) at (-2.15, -1.1) {Análise Tátil de\\Solos \& VESS};

\node[
  draw=erGoldDark, fill=bgGold, rounded corners=12pt, line width=1.15pt,
  text width=3.4cm, minimum height=1.18cm, inner sep=7pt, align=center,
  font=\footnotesize\bfseries\sffamily, text=erGoldDark
] (p3) at (2.15, -1.1) {Cadernetas\\Agroecológicas};

\node[
  draw=erForest, fill=bgForest, rounded corners=12pt, line width=1.15pt,
  text width=3.4cm, minimum height=1.18cm, inner sep=7pt, align=center,
  font=\footnotesize\bfseries\sffamily, text=erForest
] (p4) at (6.45, -1.1) {Unidades Demonstrativas\\\& Maquetes};

% Tronco vertical e barramento horizontal limpo sem sobreposição de cores
\draw[line width=1.2pt, color=erGoldDark!75] (hub.south) -- (0, 0.22);
\node[
  fill=white, text=erMuted, font=\scriptsize\bfseries\sffamily, inner xsep=8pt, inner ysep=3pt
] (exlabel) at (0, 0.72) {EXEMPLOS DE DISPOSITIVOS DE MEDIAÇÃO};

\draw[line width=1.15pt, color=erMuted!55] (-6.45, 0.22) -- (6.45, 0.22);
\draw[line width=1.15pt, color=erTeal,     ->, >=Stealth] (-6.45, 0.22) -- (p1.north);
\draw[line width=1.15pt, color=erNavy,     ->, >=Stealth] (-2.15, 0.22) -- (p2.north);
\draw[line width=1.15pt, color=erGoldDark, ->, >=Stealth] ( 2.15, 0.22) -- (p3.north);
\draw[line width=1.15pt, color=erForest,   ->, >=Stealth] ( 6.45, 0.22) -- (p4.north);

\end{tikzpicture}
\end{document}
